\chapter{Integral Lebesgue}\label{ch:b3:lebesgue}

Integral Riemann mengiris \emph{daerah asalnya} menjadi selang
kecil; sedangkan integral Lebesgue mengiris \emph{daerah hasilnya}: yakni untuk mengintegralkan
$f$, ukurlah himpunan $\{f > t\}$. Perubahannya tampak tak berbahaya
padahal revolusioner. Limit dan integral, yang selamanya bertengkar
pada teori Riemann (sebab kekonvergenan seragam dituntut!), kini
diperdamaikan oleh tiga teorema kekonvergenan --- yakni kekonvergenan
monoton, Fatou, dan kekonvergenan terdominasi --- yang hipotesisnya
lemah sampai hampir memalukan. Bab ini membangun
integralnya atas sembarang \omterm{def:b3:measure:measure}{ruang ukuran} $(X, \mathcal A,
\mu)$, membuktikan ketiga teorema itu, menuntaskan hubungannya yang
persis dengan integral Riemann (bahwa fungsi terbatas
\omterm{def:b3:lebesgue:l1}{terintegralkan} Riemann jika dan hanya jika ia \omterm{def:b3:topology:continuity}{kontinu} \omterm{def:b3:lebesgue:l1}{hampir di mana-mana}), dan
mengindustrikan pendiferensialan integral bergantung parameter
--- yakni teknik yang dipakai soal akhir pekan untuk
menghitung $\int_0^\infty\frac{\sin x}x\,\dd x$ dan
$\int_\R \eu^{-x^2}\dd x$.

\section{Fungsi terukur}

\begin{definition}\label{def:b3:lebesgue:measurable}
Misalkan $(X, \mathcal A)$ dan $(Y, \mathcal B)$ ruang terukur.
Pemetaan $f \colon X \to Y$ disebut \emph{terukur}\index{fungsi
terukur} jika $f^{-1}(B) \in \mathcal A$ untuk setiap $B \in
\mathcal B$. Untuk fungsi real (atau yang bernilai di
$[-\infty,+\infty]$), $Y = \R$ membawa aljabar-$\sigma$ Borelnya, dan
cukup diperiksa $f^{-1}(\intoo t{+\infty}) = \{f > t\} \in
\mathcal A$ untuk setiap $t \in \R$: sebab himpunan baiknya $\{B :
f^{-1}(B) \in \mathcal A\}$ membentuk aljabar-$\sigma$
(karena prapeta komutatif dengan operasi himpunan) yang memuat
sinar pembangkitnya (\cref{def:b3:measure:borel},
\cref{met:b3:measure:goodsets}).
\end{definition}

\begin{proposition}\label{prop:b3:lebesgue:stability}
(a) Komposisi pemetaan \omterm{def:b3:lebesgue:measurable}{terukur} tetap \omterm{def:b3:lebesgue:measurable}{terukur}; dan \omterm{def:b3:topology:continuity}{pemetaan
kontinu} bersifat \omterm{def:b3:lebesgue:measurable}{terukur} Borel.
(b) Jika $f, g \colon X \to \R$ \omterm{def:b3:lebesgue:measurable}{terukur}, maka demikian pula $f + g$,
$fg$, $\max(f,g)$, $\abs f$, dan $\lambda f$.
(c) Jika $(f_n)$ \omterm{def:b3:lebesgue:measurable}{terukur} dengan nilai di
$[-\infty, +\infty]$, maka $\sup_nf_n$, $\inf_nf_n$,
$\limsup f_n$, $\liminf f_n$ \omterm{def:b3:lebesgue:measurable}{terukur}; dan jika $f_n \to f$
titik demi titik, maka $f$ \omterm{def:b3:lebesgue:measurable}{terukur}.
\end{proposition}

\begin{proof}
(a) $(g\circ f)^{-1}(B) = f^{-1}(g^{-1}(B))$; sedangkan \omterm{def:b3:topology:continuity}{kekontinuannya} memberikan
\omterm{def:b3:lebesgue:measurable}{keterukuran} lewat \omterm{def:b3:topology:topology}{himpunan terbuka} pembangkitnya
(\cref{pb:b3:measure:1}, pertanyaan 10, dalam bentuk umumnya).
(b) Pemetaan $(f, g) \colon X \to \R^2$ \omterm{def:b3:lebesgue:measurable}{terukur} terhadap aljabar-$\sigma$
Borel dari $\R^2$ --- periksalah pada kotak terbukanya, yang
membangkitkannya (sebab \omterm{def:b3:topology:topology}{himpunan terbuka} $\R^2$ berupa gabungan terbilang kotak
rasional): $(f,g)^{-1}(U\times V) = f^{-1}(U)\cap g^{-1}(V)$ ---
dan $+, \times, \max$ \omterm{def:b3:topology:continuity}{kontinu} $\R^2 \to \R$: lalu komposisikan.
(c) $\{\sup f_n > t\} = \bigcup_n\{f_n > t\}$; lalu $\inf = -\sup(-)$;
$\limsup = \inf_N\sup_{n \geq N}$; dan limit titik demi titik adalah
$\limsup$-nya sendiri.
\end{proof}

\begin{definition}\label{def:b3:lebesgue:simple}
Sebuah \emph{fungsi sederhana}\index{fungsi sederhana} adalah \omterm{def:b3:lebesgue:measurable}{fungsi
terukur} yang bernilai berhingga banyak: $s = \sum_{i=1}^n
c_i\,\mathbf 1_{A_i}$, dengan $A_i \in \mathcal A$ saling lepas dan $c_i \geq
0$ (untuk teori tak negatifnya). Integralnya adalah
\[
\int s\,\dd\mu = \sum_i c_i\,\mu(A_i) \in [0, +\infty]
\]
(dengan kesepakatan $0\cdot\infty = 0$); dan nilainya tak bergantung pada
penyajiannya (haluskan kedua partisinya).
\end{definition}

\begin{theorem}[Penghampiran oleh fungsi sederhana]
\label{thm:b3:lebesgue:approximation}
Setiap $f \colon X \to [0, +\infty]$ yang \omterm{def:b3:lebesgue:measurable}{terukur} merupakan limit titik
demi titik sebuah barisan \emph{naik} berisi \omterm{def:b3:lebesgue:simple}{fungsi sederhana}:
\[
s_n = \sum_{k=1}^{n2^n} \frac{k-1}{2^n}\,
\mathbf 1_{\{\frac{k-1}{2^n} \leq f < \frac k{2^n}\}}
+ n\,\mathbf 1_{\{f \geq n\}} \nearrow f .
\]
\end{theorem}

\begin{proof}
Setiap $s_n$ bersifat sederhana (sebab himpunannya prapeta himpunan Borel).
Kemonotonannya: beralih dari $n$ ke $n+1$ membelah setiap
aras diadiknya menjadi dua dan tak pernah menurunkan nilai yang diberikan (sebab titik
dengan $\frac{k-1}{2^n} \leq f(x) < \frac k{2^n}$ memperoleh entah
$\frac{2k-2}{2^{n+1}}$ entah $\frac{2k-1}{2^{n+1}}$, keduanya $\geq
\frac{k-1}{2^n}$; dan batas atasnya $n$ juga naik). Kekonvergenannya: jika
$f(x) < \infty$, maka untuk $n > f(x)$ berlaku $f(x) - s_n(x) \leq
2^{-n}$; sedangkan jika $f(x) = \infty$, maka $s_n(x) = n \to \infty$.
\end{proof}

\section{Integralnya dan teorema kekonvergenannya}

\begin{definition}\label{def:b3:lebesgue:integral}
Untuk $f \geq 0$ yang \omterm{def:b3:lebesgue:measurable}{terukur}:
\[
\int f \,\dd\mu = \sup\Bigl\{\int s\,\dd\mu : s \text{ sederhana},
\ 0 \leq s \leq f\Bigr\} \in [0, +\infty].
\]
Integral ini monoton terhadap $f$ menurut konstruksinya, dan memperluas kasus
sederhananya (sebab untuk $f$ yang sederhana, supremumnya tercapai di $f$: bandingkan
integral sederhananya lewat penghalusan bersama).
\end{definition}

\begin{theorem}[Kekonvergenan monoton, Beppo Levi]
\label{thm:b3:lebesgue:mct}
Jika $0 \leq f_n \nearrow f$ titik demi titik (dan \omterm{def:b3:lebesgue:measurable}{terukur}), maka
\[
\int f_n\,\dd\mu \nearrow \int f\,\dd\mu .
\]
\index{teorema kekonvergenan monoton}
\end{theorem}

\begin{proof}
Fungsi $f$ \omterm{def:b3:lebesgue:measurable}{terukur} (\cref{prop:b3:lebesgue:stability}(c)) dan
$\int f_n$ naik menuju suatu $L \leq \int f$ (menurut kemonotonannya).
Sebaliknya, tetapkan sebuah $s = \sum c_i\mathbf 1_{A_i} \leq f$ yang sederhana
dan $\theta \in (0,1)$; maka himpunan $E_n = \{f_n \geq \theta s\}$
\omterm{def:b3:lebesgue:measurable}{terukur} dan naik menuju $X$ (sebab bila $s(x) > 0$: $f(x)
\geq s(x) > \theta s(x)$, sehingga akhirnya $f_n(x) \geq \theta
s(x)$; sedangkan bila $s(x) = 0$: secara trivial). Maka
\[
\int f_n \geq \int_{E_n}\theta s\,\dd\mu
= \theta\sum_i c_i\,\mu(A_i \cap E_n)
\xrightarrow[n\to\infty]{} \theta\sum_ic_i\,\mu(A_i)
= \theta\int s
\]
menurut \omterm{def:b3:topology:continuity}{kekontinuan} dari bawah (\cref{prop:b3:measure:basics}(c)).
Jadi $L \geq \theta\int s$ untuk setiap $\theta < 1$ dan setiap $s
\leq f$ yang sederhana: sehingga $L \geq \int f$.
\end{proof}

\begin{corollary}\label{cor:b3:lebesgue:additivity}
Untuk $f, g \geq 0$ yang \omterm{def:b3:lebesgue:measurable}{terukur} dan $c \geq 0$: berlaku $\int(f + g) =
\int f + \int g$ dan $\int cf = c\int f$; sedangkan untuk deret
\omterm{def:b3:lebesgue:measurable}{fungsi terukur} tak negatif, $\int\sum_nf_n =
\sum_n\int f_n$.
\end{corollary}

\begin{proof}
Untuk \omterm{def:b3:lebesgue:simple}{fungsi sederhana}, keaditifannya berupa perhitungan pada
penghalusan bersama. Secara umum ambillah $s_n \nearrow f$ dan $t_n \nearrow g$
(\cref{thm:b3:lebesgue:approximation}): maka $s_n + t_n \nearrow f +
g$, dan teorema kekonvergenan monotonnya mengalihkan keaditifannya ke limitnya. Pernyataan
deretnya adalah teorema itu yang diterapkan pada jumlah parsialnya.
\end{proof}

\begin{theorem}[Lema Fatou]\label{thm:b3:lebesgue:fatou}
Untuk $f_n \geq 0$ yang \omterm{def:b3:lebesgue:measurable}{terukur}:
\[
\int \liminf_n f_n \,\dd\mu \;\leq\; \liminf_n \int
f_n\,\dd\mu .
\]
\index{lema Fatou}
\end{theorem}

\begin{proof}
Misalkan $g_N = \inf_{n\geq N}f_n$: yang \omterm{def:b3:lebesgue:measurable}{terukur}, dengan $0 \leq g_N \nearrow
\liminf f_n$, dan $g_N \leq f_n$ untuk setiap $n \geq N$, sehingga
$\int g_N \leq \inf_{n \geq N}\int f_n$. Terapkan kekonvergenan monoton pada ruas
kirinya: $\int\liminf f_n = \lim_N\int g_N \leq
\lim_N\inf_{n\geq N}\int f_n = \liminf\int f_n$.
\end{proof}

\begin{definition}\label{def:b3:lebesgue:l1}
Sebuah $f \colon X \to \R$ (atau $\C$) yang \omterm{def:b3:lebesgue:measurable}{terukur} disebut
\emph{terintegralkan}\index{fungsi terintegralkan} jika $\int\abs
f\,\dd\mu < \infty$; lalu $\int f = \int f^+ - \int f^-$
(yakni bagian positif dan negatifnya; atau bagian real dan imajinernya pada
kasus kompleksnya). Integralnya bersifat linear pada fungsi terintegralkan
(uraikan lalu gabungkan kembali bagian positifnya; sedangkan kasus kompleksnya
menyusut ke kasus realnya) dan memenuhi $\abs{\int f} \leq
\int\abs f$ (pada kasus realnya: $\pm\int f = \int(\pm f) \leq
\int\abs f$; sedangkan pada kasus kompleksnya: kalikan dengan konstanta unimodular
agar integralnya menjadi real). Sebuah sifat berlaku \emph{hampir
di mana-mana}\index{hampir di mana-mana} jika ia gagal hanya pada
himpunan nol-$\mu$; dan mengubah $f$ pada himpunan nol tak mengubah
integral apa pun (sebab selisihnya didominasi $\infty\cdot\mathbf
1_N$, yang berintegral $0$).
\end{definition}

\begin{theorem}[Kekonvergenan terdominasi]\label{thm:b3:lebesgue:dct}
Misalkan $f_n \to f$ \omterm{def:b3:lebesgue:l1}{hampir di mana-mana}, dengan $\abs{f_n} \leq g$ \omterm{def:b3:lebesgue:l1}{hampir di mana-mana} untuk sebuah
$g$ \emph{\omterm{def:b3:lebesgue:l1}{terintegralkan}} yang tetap. Maka $f$ \omterm{def:b3:lebesgue:l1}{terintegralkan} dan
\[
\int f_n\,\dd\mu \longrightarrow \int f\,\dd\mu,
\qquad\text{bahkan}\quad \int\abs{f_n - f}\,\dd\mu \to 0 .
\]
\index{teorema kekonvergenan terdominasi}
\end{theorem}

\begin{proof}
Buanglah sebuah himpunan nol agar hipotesisnya berlaku titik demi titik. Dari $\abs f
\leq g$: $f$ \omterm{def:b3:lebesgue:l1}{terintegralkan}. Fungsi $h_n = 2g - \abs{f_n
- f} \geq 0$ memenuhi $\liminf h_n = 2g$; lalu Fatou memberikan
\[
\int 2g \leq \liminf\int\bigl(2g - \abs{f_n - f}\bigr)
= \int 2g - \limsup\int\abs{f_n - f},
\]
sehingga $\limsup\int\abs{f_n - f} \leq 0$ (dan pengurangannya sah sebab
$\int 2g < \infty$). Akhirnya $\abs{\int f_n - \int f} \leq
\int\abs{f_n - f} \to 0$.
\end{proof}

\begin{method}\label{met:b3:lebesgue:convergence}
Menghadapi $\lim_n\int f_n$: cobalah, berurutan --- (1) apakah
barisannya monoton (atau berupa deret bersuku tak negatif)? Pakai kekonvergenan monoton, tanpa perlu
keterintegralan. (2) Adakah satu pendominasi tunggal yang \omterm{def:b3:lebesgue:l1}{terintegralkan}
$g \geq \abs{f_n}$, yang ditemukan lewat batas kasar
(yakni ``$\sup_n$'' atas taksirannya)? Pakai kekonvergenan terdominasi. (3) Tanpa dominasi dan tanpa
kemonotonan? Fatou tetap membatasi salah satu sisinya, dan kesamaannya dapat
sungguh gagal: sebab gundukan yang lolos $f_n = n\mathbf
1_{\intoo0{1/n}}$ mempunyai $\int f_n = 1$ padahal $f_n \to 0$ \omterm{def:b3:lebesgue:l1}{hampir di mana-mana}.
Jadi dominasi tepat merupakan hal yang melarang massanya lolos ke
tak hingga, yang terjadi secara tegak ataupun mendatar.
\end{method}

\section{Riemann lawan Lebesgue}

\begin{theorem}[Kriteria Lebesgue]
\label{thm:b3:lebesgue:riemann}
Misalkan $f \colon \intcc ab \to \R$ \emph{terbatas}. Maka $f$
\omterm{def:b3:lebesgue:l1}{terintegralkan} Riemann jika dan hanya jika $f$ \omterm{def:b3:topology:continuity}{kontinu} \omterm{def:b3:lebesgue:l1}{hampir di mana-mana}
terhadap $\lambda$; dan dalam hal itu $f$ \omterm{def:b3:lebesgue:l1}{terintegralkan} Lebesgue serta kedua
integralnya berimpit.\index{integral Riemann lawan Lebesgue}
\end{theorem}

\begin{proof}
Untuk sebuah pembagian $\sigma = (a = x_0 < \dots < x_N = b)$, misalkan
$L_\sigma$ dan $U_\sigma$ fungsi tangga yang bernilai, pada setiap
$\intoo{x_{i-1}}{x_i}$, sebesar $m_i = \inf_{[x_{i-1}, x_i]}f$ dan
$M_i = \sup$; maka jumlah Darbouxnya adalah integralnya (sebab Riemann dan
Lebesgue bersesuaian pada fungsi tangga, keduanya memberikan $\sum
m_i\Delta x_i$). Ambil sebuah barisan pembagian $\sigma_n$,
yang masing-masing menghaluskan pendahulunya, dengan kehalusan $\to 0$, sehingga jumlah Darbouxnya
konvergen ke integral Darboux bawah dan atas dari $f$.
Penghalusannya membuat $L_{\sigma_n}$ tak turun dan
$U_{\sigma_n}$ tak naik secara titik demi titik di luar himpunan terbilang
$D$ berisi semua titik pembaginya; sebutlah limitnya $\ell$ dan $u$
(yang \omterm{def:b3:lebesgue:measurable}{terukur}, \cref{prop:b3:lebesgue:stability}). Untuk $x \notin
D$, dengan menulis $I_n(x)$ untuk selang-$\sigma_n$ terbuka yang
memuat $x$: berlaku $\ell(x) = \sup_n\inf_{I_n(x)}f$ dan $u(x) =
\inf_n\sup_{I_n(x)}f$; dan karena kehalusannya menyusut ke $0$, keduanya
merupakan \emph{selubung} bawah dan atas dari $f$ di $x$
--- sehingga $u(x) - \ell(x)$ adalah osilasi $f$ di $x$ --- sehingga
\emph{$\ell(x) = u(x)$ jika dan hanya jika $f$ \omterm{def:b3:topology:continuity}{kontinu} di $x$}. Lalu lewat kekonvergenan monoton/terdominasi (karena terbatas, pada selang berhingga):
\[
\int_{\intcc ab}\ell\,\dd\lambda = \lim_n\int L_{\sigma_n}
= \underline{\int}f,
\qquad
\int_{\intcc ab}u\,\dd\lambda = \overline{\int}f .
\]
Jadi $f$ \omterm{def:b3:lebesgue:l1}{terintegralkan} Riemann $\iff$ $\underline\int f =
\overline\int f$ $\iff$ $\int(u - \ell) = 0$ $\iff$ $u = \ell$
\omterm{def:b3:lebesgue:l1}{hampir di mana-mana} (sebab $u - \ell \geq 0$; \cref{exo:b3:lebesgue:5}) $\iff$ $f$
\omterm{def:b3:topology:continuity}{kontinu} \omterm{def:b3:lebesgue:l1}{hampir di mana-mana}. Dalam hal itu $\ell \leq f \leq u$ dengan $\ell =
u$ \omterm{def:b3:lebesgue:l1}{hampir di mana-mana}: sehingga $f$ sama dengan $\ell$ yang \omterm{def:b3:lebesgue:measurable}{terukur} \omterm{def:b3:lebesgue:l1}{hampir di mana-mana}, jadi
\omterm{def:b3:lebesgue:measurable}{terukur} Lebesgue (menurut \omterm{def:b3:complete:complete}{kelengkapan} $\lambda$) dengan
$\int f\,\dd\lambda = \int\ell\,\dd\lambda = \underline\int f =
\int_a^bf$.
\end{proof}

\begin{example}\label{ex:b3:lebesgue:riemannexamples}
Fungsi $\mathbf 1_\Q$ tak \omterm{def:b3:topology:continuity}{kontinu} di mana pun: jadi tak \omterm{def:b3:lebesgue:l1}{terintegralkan} Riemann
--- tetapi trivial bagi Lebesgue: $\int\mathbf 1_\Q\,\dd\lambda =
\lambda(\Q) = 0$. Fungsi Thomae (yakni $\frac1q$ di bilangan rasional
$\frac pq$ dan $0$ selainnya) \omterm{def:b3:topology:continuity}{kontinu} tepat di
bilangan irasional: sehingga \omterm{def:b3:lebesgue:l1}{terintegralkan} Riemann dengan integral $0$. Dan
integral Riemann \emph{tak wajar} merupakan gagasan yang berbeda:
$\int_0^\infty\frac{\sin x}x\,\dd x$ konvergen sebagai limit
$\int_0^A$ (dan soal akhir pekan menghitungnya $= \frac\pi2$),
padahal $\frac{\sin x}x \notin L^1(\intoo0{+\infty})$: sebab integral
mutlaknya divergen seperti deret harmonik
(\cref{exo:b3:lebesgue:6}). Jadi teori Lebesgue menukar
kekonvergenan bersyarat dengan teorema limit yang tegar.
\end{example}

\section{Integral berparameter}

Sepanjang bagian ini, $(X, \mathcal A, \mu)$ adalah \omterm{def:b3:measure:measure}{ruang ukuran}, $T$ sebuah
ruang metrik (yakni parameternya), dan $f \colon T \times X \to \C$
dengan $f(t, \cdot)$ \omterm{def:b3:lebesgue:l1}{terintegralkan} untuk setiap $t$; lalu tetapkan $F(t) = \int_X
f(t, x)\,\dd\mu(x)$.

\begin{theorem}[Kekontinuan]\label{thm:b3:lebesgue:paramcont}
Andaikan: $t \mapsto f(t,x)$ \omterm{def:b3:topology:continuity}{kontinu} di $t_0$ untuk hampir setiap
$x$, dan ada $g$ \omterm{def:b3:lebesgue:l1}{terintegralkan} dengan $\abs{f(t,x)} \leq
g(x)$ untuk setiap $t$ pada sebuah \omterm{def:b3:topology:topology}{persekitaran} $t_0$ dan hampir setiap $x$.
Maka $F$ \omterm{def:b3:topology:continuity}{kontinu} di $t_0$.\index{integral berparameter}
\end{theorem}

\begin{proof}
Untuk sembarang barisan $t_n \to t_0$: $f(t_n, \cdot) \to f(t_0,
\cdot)$ \omterm{def:b3:lebesgue:l1}{hampir di mana-mana}, dan terdominasi oleh $g$: sehingga kekonvergenan terdominasinya memberikan $F(t_n) \to F(t_0)$;
dan \omterm{def:b3:topology:continuity}{kekontinuan} barisan sudah cukup di ruang metrik
(\cref{rem:b3:topology:sequences}).
\end{proof}

\begin{theorem}[Pendiferensialan di bawah integral]
\label{thm:b3:lebesgue:paramdiff}
Misalkan $T$ sebuah selang terbuka $\R$. Andaikan: untuk hampir setiap $x$,
$t \mapsto f(t,x)$ dapat diturunkan pada $T$, dengan
\[
\Bigl|\frac{\partial f}{\partial t}(t, x)\Bigr| \leq g(x)
\quad \text{untuk setiap } t \in T,\ \text{hampir setiap } x,
\]
dengan $g$ \omterm{def:b3:lebesgue:l1}{terintegralkan}. Maka $F$ dapat diturunkan pada $T$ dengan $F'(t)
= \int_X \frac{\partial f}{\partial t}(t, x)\,\dd\mu(x)$.
\end{theorem}

\begin{proof}
Tetapkan $t$ dan $h_n \to 0$: maka hasil bagi selisihnya
\[
\varphi_n(x) = \frac{f(t + h_n, x) - f(t, x)}{h_n}
\longrightarrow \frac{\partial f}{\partial t}(t,x)
\quad\text{hampir di mana-mana},
\]
dan ketaksamaan nilai rata-rata membatasi $\abs{\varphi_n(x)} \leq
\sup_{s}\abs{\partial_tf(s,x)} \leq g(x)$: sehingga kekonvergenan terdominasinya berlaku, dan
$\frac{F(t + h_n) - F(t)}{h_n} = \int\varphi_n \to
\int\partial_t f(t, \cdot)$.
\end{proof}

\begin{example}[Fungsi Gamma]\label{ex:b3:lebesgue:gamma}
Untuk $t > 0$ misalkan\index{fungsi Gamma}
\[
\Gamma(t) = \int_0^{+\infty} x^{t-1}\eu^{-x}\,\dd x .
\]
Integralnya konvergen: sebab di dekat $0$, $x^{t-1}$ \omterm{def:b3:lebesgue:l1}{terintegralkan} (karena $t >
0$); sedangkan di tak hingga, $x^{t-1}\eu^{-x} \leq C\eu^{-x/2}$.
Pengintegralan parsial (pada $[\varepsilon, A]$, lalu limitnya lewat
kekonvergenan monoton) memberikan persamaan fungsional $\Gamma(t + 1) =
t\,\Gamma(t)$, sehingga $\Gamma(n+1) = n!$: yakni faktorial yang
terinterpolasi. Pada setiap $\intcc ab \subseteq \intoo0{+\infty}$,
$\partial_t\bigl(x^{t-1}\eu^{-x}\bigr) = \ln
x\cdot x^{t-1}\eu^{-x}$ terdominasi oleh $\abs{\ln x}(x^{a-1} +
x^{b-1})\eu^{-x}$ yang \omterm{def:b3:lebesgue:l1}{terintegralkan}: sehingga $\Gamma$ bersifat $\mathcal C^1$, dan
lewat induksi $\mathcal C^\infty$, dengan $\Gamma^{(k)}(t) =
\int_0^\infty(\ln x)^kx^{t-1}\eu^{-x}\dd x$. Sedangkan nilai
$\Gamma(\frac12) = \sqrt\pi$ merupakan integral Gauss yang
menyamar (\cref{pb:b3:lebesgue:1}).
\end{example}

\begin{omfigure}
\begin{tikzpicture}
\begin{axis}[omaxis, width=10.6cm, height=5.4cm,
    xmin=0, xmax=25, ymin=-0.25, ymax=1.05,
    xtick={0, 3.1416, 6.2832, 9.4248, 12.566, 15.708, 18.850,
      21.991, 25.133},
    xticklabels={$0$, $\pi$, $2\pi$, $3\pi$, $4\pi$, $5\pi$,
      $6\pi$, $7\pi$, $8\pi$},
    ytick={0, 1}]
  \addplot[omDef!40] coordinates {(0,0) (25,0)};
  \addplot[omThm, thick, domain=0.02:25, samples=600]
    {sin(deg(x))/x};
  \addplot[omDef, dashed, domain=0.02:25, samples=200]
    {1/x};
  \addplot[omDef, dashed, domain=0.02:25, samples=200]
    {-1/x};
\end{axis}
\end{tikzpicture}

{\small Integran $\frac{\sin x}x$: integral tak wajarnya
$\int_0^\infty$ konvergen lewat peniadaan berganti tanda antara
lengkungannya, tetapi luas $\abs{\cdot}$ setiap lengkungannya berperilaku
seperti $\frac2{\pi k}$ --- yakni deret harmonik: sehingga $\frac{\sin x}x
\notin L^1$. Jadi keterintegralan Lebesgue adalah keterintegralan mutlak.}
\end{omfigure}

\section{Latihan}

\begin{exercise}[$\star$]\label{exo:b3:lebesgue:1}
(a) Tunjukkan bahwa fungsi monoton $\R \to \R$ bersifat
\omterm{def:b3:lebesgue:measurable}{terukur} Borel, dan bahwa turunan (dari fungsi yang dapat diturunkan
di mana-mana) bersifat \omterm{def:b3:lebesgue:measurable}{terukur} Borel.
(b) Tunjukkan bahwa $f \colon X \to \R$ \omterm{def:b3:lebesgue:measurable}{terukur} jika dan hanya jika $\{f > q\}
\in \mathcal A$ untuk setiap $q$ \emph{rasional}.
\end{exercise}

\begin{exercise}[$\star$]\label{exo:b3:lebesgue:2}
Hitunglah, dengan pembenaran penuh:
\[
\lim_{n\to\infty}\int_0^{+\infty}
\frac{\cos x}{(1 + x/n)^{n}}\,\dd x,
\qquad
\lim_{n\to\infty}\int_0^1 \frac{n\,x^{n-1}}{1 + x}\,\dd x .
\]
\emph{(Untuk yang kedua: substitusikan $u = x^n$ sebelum mendominasinya.)}
\end{exercise}

\begin{exercise}[$\star\star$]\label{exo:b3:lebesgue:3}
(a) Tampilkan ketaksamaan sejati pada lema Fatou.
(b) Tampilkan $f_n \to 0$ titik demi titik dengan $\int f_n = 1$ dalam tiga
cara: yakni lolos secara tinggi, secara lebar, dan ke tak hingga. Hipotesis
tunggal mana pada kekonvergenan terdominasi yang dilanggar masing-masingnya?
(c) Tunjukkan bahwa pada lema Fatou kita tak dapat mengganti $\liminf$ dengan
$\limsup$ pada ruas mana pun.
\end{exercise}

\begin{exercise}[$\star\star$]\label{exo:b3:lebesgue:4}
(a) Tunjukkan bahwa $\displaystyle\int_0^{+\infty}\frac{x}{\eu^x -
1}\,\dd x = \sum_{n\geq1}\frac1{n^2} = \frac{\pi^2}6$
\emph{(uraikan $\frac1{\eu^x - 1}$ menjadi deret geometri lalu
integralkan suku demi suku --- teorema mana yang mengizinkannya?)}.
(b) (Mimpi mahasiswa tahun kedua) Tunjukkan bahwa
$\displaystyle\int_0^1 x^{-x}\,\dd x = \sum_{n\geq1}n^{-n}$.
\emph{(Tulislah $x^{-x} = \eu^{-x\ln x} = \sum_k\frac{(-x\ln
x)^k}{k!}$ lalu hitunglah $\int_0^1(-x\ln x)^k\dd x$ dengan
substitusi $x = \eu^{-u/(k+1)}$, sambil mengenali $\Gamma$.)}
\end{exercise}

\begin{exercise}[$\star\star$]\label{exo:b3:lebesgue:5}
(a) Tunjukkan bahwa $f \geq 0$ yang \omterm{def:b3:lebesgue:measurable}{terukur} dengan $\int f\,\dd\mu = 0$
memenuhi $f = 0$ \omterm{def:b3:lebesgue:l1}{hampir di mana-mana}. \emph{(Tinjaulah $\{f \geq 1/n\}$ dan
ketaksamaan Markov: $\mu(\{f \geq a\}) \leq \frac1a\int f$ ---
lalu buktikanlah.)}
(b) Tunjukkan bahwa $f$ yang \omterm{def:b3:lebesgue:l1}{terintegralkan} bernilai berhingga \omterm{def:b3:lebesgue:l1}{hampir di mana-mana}.
(c) Tunjukkan bahwa jika $\int_A f\,\dd\mu = 0$ untuk \emph{setiap}
$A$ yang \omterm{def:b3:lebesgue:measurable}{terukur}, maka $f = 0$ \omterm{def:b3:lebesgue:l1}{hampir di mana-mana}.
\end{exercise}

\begin{exercise}[$\star\star$]\label{exo:b3:lebesgue:6}
(a) Terapkan \cref{thm:b3:lebesgue:riemann} untuk memutuskan keterintegralan
Riemann dari: $\mathbf 1_\Q$; fungsi Thomae; dan $\mathbf
1_K$ dengan $K$ sebuah himpunan Cantor gemuk (\cref{exo:b3:measure:5}).
(b) Tunjukkan bahwa $\int_1^{+\infty}\abs{\frac{\sin x}x}\,\dd x =
+\infty$, sedangkan $\lim_{A\to\infty}\int_1^A\frac{\sin
x}x\,\dd x$ ada \emph{(integralkan secara parsial)}: yakni kekonvergenan
tak wajar tanpa keterintegralan.
\end{exercise}

\begin{exercise}[$\star\star$]\label{exo:b3:lebesgue:7}
Benarkanlah bahwa $F(t) = \int_0^{+\infty}\eu^{-x^2}\cos(tx)\,\dd x$
bersifat $\mathcal C^1$ pada $\R$ dan memenuhi $F'(t) = -\frac t2F(t)$
\emph{(integralkan secara parsial)}; lalu simpulkan $F(t) =
F(0)\,\eu^{-t^2/4}$. (Dengan $F(0) = \frac{\sqrt\pi}2$ dari
soal akhir pekannya: fungsi Gauss pada dasarnya merupakan transformasi
Fouriernya sendiri --- dan \cref{ch:b3:fouriertransform} akan
menyistematiskannya.)
\end{exercise}

\begin{exercise}[$\star\star\star$]\label{exo:b3:lebesgue:8}
(Frullani) Misalkan $0 < a < b$. Tunjukkan bahwa
\[
\int_0^{+\infty}\frac{\eu^{-ax} - \eu^{-bx}}{x}\,\dd x =
\ln\frac ba,
\]
dengan menulis integrannya sebagai $\int_a^b \eu^{-xt}\,\dd t$ lalu
membenarkan penukarannya lewat teori tak negatifnya
(\cref{cor:b3:lebesgue:additivity} dalam bentuk \omterm{def:b3:topology:continuity}{kontinu} ---
dahuluilah Tonelli, atau irislah $[a,b]$ menjadi $n$ bagian sama besar lalu
ambil limitnya).
\end{exercise}

\begin{exercise}[$\star\star$]\label{exo:b3:lebesgue:9}
Misalkan $f \geq 0$ \omterm{def:b3:lebesgue:measurable}{terukur} pada $(X, \mathcal A, \mu)$. Tunjukkan
bahwa $\nu(A) = \int_A f\,\dd\mu$ mendefinisikan sebuah \omterm{def:b3:measure:measure}{ukuran}
(dengan \emph{kepadatan}\index{kepadatan ukuran} $f$ terhadap
$\mu$), dan bahwa $\int g\,\dd\nu = \int gf\,\dd\mu$ untuk setiap
$g \geq 0$ yang \omterm{def:b3:lebesgue:measurable}{terukur} \emph{(buktikanlah untuk fungsi indikator, lalu
\omterm{def:b3:lebesgue:simple}{fungsi sederhana}, lalu pakai kekonvergenan monoton --- yakni mesin bakunya)}.
\end{exercise}

\begin{exercise}[$\star\star\star$]\label{exo:b3:lebesgue:10}
(Sebuah kegagalan bergaya Weierstrass) Definisikan $f(t) =
\int_0^{+\infty}\frac{\sin(tx)}{x(1 + x^2)}\,\dd x$.
(a) Tunjukkan bahwa $f$ terdefinisi dengan baik dan \omterm{def:b3:topology:continuity}{kontinu} pada $\R$, serta
bersifat $\mathcal C^1$ dengan $f'(t) = \int_0^\infty\frac{\cos(tx)}{1 +
x^2}\dd x$ untuk setiap $t$ --- tetapi bahwa menurunkannya
\emph{sekali lagi} di bawah integralnya tidak sah.
(b) Dengan menerima $f'(t) = \frac\pi2\eu^{-t}$ untuk $t > 0$ (yang dibuktikan
di \cref{ch:b3:residues}), berapakah
$\int_0^\infty\frac{x\sin(tx)}{1+x^2}\dd x$ untuk $t > 0$, dan
mengapa rumusnya membenarkan kegagalan pada (a)?
\end{exercise}

\begin{exercise}[$\star\star$]\label{exo:b3:lebesgue:11}
(Lema Scheffé) Misalkan $f_n, f \geq 0$ \omterm{def:b3:lebesgue:l1}{terintegralkan} dengan
$f_n \to f$ \omterm{def:b3:lebesgue:l1}{hampir di mana-mana} dan $\int f_n \to \int f$.
(a) Tunjukkan bahwa $\int\abs{f_n - f} \to 0$. \emph{(Terapkan
kekonvergenan terdominasi pada $g_n = (f - f_n)^+ \leq f$, lalu
tulislah $\int\abs{f_n - f} = 2\int g_n - \int(f - f_n)$.)}
(b) Tunjukkan lewat contoh bahwa hipotesis $\int f_n \to \int
f$ tak dapat dibuang (yakni gundukan yang meluncur atau memusat), dan
bahwa kesimpulannya gagal bagi $f_n$ bertanda tanpa
kendali nilai mutlak: sebab $f_n = n\mathbf 1_{\intoc0{1/n}} -
n\mathbf 1_{\intoc{-1/n}0}$ mempunyai $f_n \to 0$ \omterm{def:b3:lebesgue:l1}{hampir di mana-mana}, $\int f_n
= 0 \to 0$, padahal $\int\abs{f_n} = 2$.
(c) Terapannya (pada \omterm{ex:b3:lebesgue:gamma}{kepadatan}): jika \omterm{ex:b3:lebesgue:gamma}{kepadatan} peluang $p_n
\to p$ \omterm{def:b3:lebesgue:l1}{hampir di mana-mana}, maka dengan sendirinya $\int\abs{p_n - p} \to 0$:
yakni kekonvergenan titik demi titik \omterm{ex:b3:lebesgue:gamma}{kepadatannya} adalah kekonvergenan $L^1$ ---
sebuah peningkatan kekonvergenan secara cuma-cuma.
\end{exercise}

\begin{exercise}[$\star\star$]\label{exo:b3:lebesgue:12}
Limit klasik, dengan pembenaran penuh lewat kekonvergenan monoton/terdominasi:
\[
\text{(a)}\ \lim_{n\to\infty}\int_0^n\Bigl(1 -
\frac xn\Bigr)^n\eu^{x/2}\,\dd x, \qquad
\text{(b)}\ \lim_{n\to\infty}\int_0^1\frac{n\,x^{n-1}}{1 +
x}\,\dd x,
\]
\[
\text{(c)}\ \lim_{n\to\infty}\int_0^\infty
\frac{\dd x}{(1 + x/n)^n\,x^{1/n}} .
\]
\emph{(Untuk (a): $(1 - x/n)^n \nearrow \eu^{-x}$ untuk $x$ yang
tetap --- buktikan kemonotonannya lewat $\log$; untuk (b),
integralkan secara parsial atau substitusikan $x = u^{1/n}$ lalu kenali
pemusatan di ujungnya; sedangkan untuk (c), carilah pendominasi
\omterm{def:b3:lebesgue:l1}{terintegralkan} yang berlaku untuk setiap $n \geq 2$ dengan memecahnya di $x =
1$.)}
\end{exercise}

\section{Soal: dua integral yang termasyhur}

\begin{problem}[{Soal akhir pekan --- integral Gauss dan
integral Dirichlet, lewat parameter belaka}]
\label{pb:b3:lebesgue:1}
Dua integral menguasai analisis terapan:
\[
G = \int_{-\infty}^{+\infty}\eu^{-x^2}\dd x = \sqrt\pi,
\qquad
D = \int_0^{+\infty}\frac{\sin x}{x}\,\dd x = \frac\pi2
\]
(dengan yang kedua sebagai integral tak wajar,
\cref{ex:b3:lebesgue:riemannexamples}). Keduanya kita buktikan memakai
perkakas bab ini saja.

\medskip
\noindent\textbf{Bagian I --- Fungsi Gauss.} Untuk $t \geq 0$ tetapkan
\[
A(t) = \Bigl(\int_0^t\eu^{-x^2}\dd x\Bigr)^{2},
\qquad
B(t) = \int_0^1\frac{\eu^{-t^2(1 + x^2)}}{1 + x^2}\,\dd x .
\]
\begin{enumerate}
  \item Benarkanlah bahwa $A$ dan $B$ bersifat $\mathcal C^1$ pada
        $\intoo0{+\infty}$ lalu hitunglah $A'$ dan $B'$; lalu tunjukkan
        $A'(t) + B'(t) = 0$. \emph{(Pada $B'$, substitusikan $u =
        tx$.)}
  \item Hitunglah $A(0) + B(0)$ dan $\lim_{t\to+\infty}(A + B)(t)$
        --- lalu benarkan limit di bawah integralnya pada $B$.
  \item Simpulkan $\int_0^\infty \eu^{-x^2}\dd x =
        \frac{\sqrt\pi}2$, sehingga $G = \sqrt\pi$, lalu turunkan
        $\Gamma(\tfrac12) = \sqrt\pi$ \emph{(substitusikan $x =
        u^2$ pada $\Gamma(\frac12)$)}.
\end{enumerate}

\medskip
\noindent\textbf{Bagian II --- Integral Dirichlet.} Untuk $t
\geq 0$ tetapkan
\[
F(t) = \int_0^{+\infty}\eu^{-tx}\,\frac{\sin x}{x}\,\dd x .
\]
\begin{enumerate}[resume]
  \item Tunjukkan bahwa integral pendefinisi $F(t)$ konvergen untuk
        setiap $t > 0$ sebagai integral Lebesgue, dan untuk $t = 0$
        sebagai integral tak wajar; lalu tunjukkan bahwa $D =
        \lim_{A\to\infty}\int_0^A\frac{\sin x}x\dd x$ ada
        \emph{(integralkan secara parsial pada $[\pi, A]$)}.
  \item Tunjukkan bahwa $F$ bersifat $\mathcal C^1$ pada
        $\intoo0{+\infty}$ dengan
        \[
        F'(t) = -\int_0^{+\infty}\eu^{-tx}\sin x\,\dd x
        = -\frac{1}{1 + t^2}
        \]
        \emph{(dengan dominasi pada $[t_0, \infty)$ untuk setiap $t_0 >
        0$; sedangkan integral terakhirnya lewat dua kali pengintegralan parsial atau
        lewat eksponensial kompleks)}.
  \item Tunjukkan bahwa $F(t) \to 0$ saat $t \to +\infty$, lalu turunkan $F(t)
        = \frac\pi2 - \arctan t$ pada $\intoo0{+\infty}$.
  \item Titik pelik: $D = \lim_{t\to0^+}F(t)$. Buktikanlah
        lewat kendali seragam atas ekornya: yakni untuk $0 \leq t \leq 1$
        dan $A \geq \pi$, integralkan secara parsial untuk menunjukkan bahwa
        \[
        \Bigl|\int_A^{+\infty}\eu^{-tx}\frac{\sin
        x}x\,\dd x\Bigr| \leq \frac{C}{A}
        \]
        dengan $C$ yang tak bergantung pada $t$ \emph{(turunkan
        $\frac{\eu^{-tx}}x$ lalu batasi $\abs{\cos}$ oleh $1$;
        perhatikan $t\eu^{-tx} \leq 1/x\cdot(tx\eu^{-tx})$ dengan
        $\sup_{u\geq0}u\eu^{-u} < 1$)}; lalu pecahlah $F(t) - D$
        menjadi $[0, A]$ (tempat kekonvergenan terdominasinya berlaku saat $t \to 0$) dan
        $[A, \infty)$.
  \item Simpulkan: $D = \frac\pi2$.
\end{enumerate}

\medskip
\noindent\textbf{Bagian III --- Buah hasilnya.}
\begin{enumerate}[resume]
  \item Hitunglah $\int_0^{+\infty}\frac{\sin^2x}{x^2}\,\dd x$
        \emph{(integralkan secara parsial lalu susutkan ke $D$ lewat $\sin
        2x = 2\sin x\cos x$)}.
  \item Hitunglah $\int_0^{+\infty}\frac{1 -
        \cos x}{x^2}\,\dd x$, lalu periksa kekonsistenan
        kedua hasilnya.
  \item Untuk $a > 0$, hitunglah
        $\int_0^{+\infty}\frac{\sin(ax)}x\dd x$ dan
        $\int_{-\infty}^{+\infty}\eu^{-ax^2}\dd x$, lalu catat
        aturan penskalaannya (yang akan menjadi kuda beban
        \cref{ch:b3:fouriertransform}).
  \item Terangkan dengan tepat mengapa $D$ tak dapat ditangani
        langsung oleh kekonvergenan terdominasi di $t = 0$ (sebab tak ada pendominasi \omterm{def:b3:lebesgue:l1}{terintegralkan} pada
        $[0,1]\times[0,\infty)$), dan mengapa pemecahan ekor pada
        pertanyaan 7 merupakan penggantinya yang jujur --- sebab pola ini
        (yakni ``keterintegralan seragam atas ekornya'') berulang di sepanjang
        analisis.
\end{enumerate}

\medskip
\noindent\textbf{Bagian IV --- \omterm{ex:b3:lebesgue:gamma}{Fungsi Gamma} menurut
Bohr dan Mollerup.} Fungsi $\Gamma$
(\cref{ex:b3:lebesgue:gamma}) memenuhi $\Gamma(1) = 1$ dan
$\Gamma(x+1) = x\Gamma(x)$ --- tetapi demikian pula tak berhingga banyak
fungsi lain (kalikan saja dengan goyangan berperiode $1$ apa pun). Satu syarat
kecembungan memakukan $\Gamma$ secara tunggal, dan
identitasnya yang lebih dalam lalu jatuh secara mekanis. Sebuah fungsi positif $f$
pada sebuah selang disebut \emph{log-cembung} jika $\log f$ cembung.
\begin{enumerate}[resume]
  \item Tunjukkan bahwa log-cembung mengakibatkan cembung, bahwa hasil kali
        fungsi log-cembung dan komposisinya dengan pemetaan
        afin bersifat log-cembung, dan --- lewat
        ketaksamaan Hölder dua fungsi $\int\abs{uv} \leq
        \bigl(\int\abs u^p\bigr)^{1/p}\bigl(\int\abs
        v^q\bigr)^{1/q}$, yang dibuktikan langsung dari ketaksamaan
        Young --- bahwa $\Gamma$ bersifat log-cembung pada
        $\intoo0\infty$.
  \item (Lema kemiringan) Misalkan $g$ cembung pada $\intoo0\infty$
        dengan $g(n+1) - g(n) = \log n$ untuk setiap bilangan bulat $n
        \geq 1$. Untuk $x \in \intoc01$ dan $n \geq 2$, bandingkanlah
        kemiringan $g$ sepanjang $[n-1, n]$, $[n, n+x]$ dan
        $[n, n+1]$, lalu turunkan
        \[
        x\log(n-1) \;\leq\; g(n + x) - g(n) \;\leq\; x\log n .
        \]
  \item (Bohr--Mollerup) Misalkan $f > 0$ memenuhi $f(1) = 1$,
        $f(x+1) = xf(x)$, dan $\log f$ cembung. Dengan mengurai
        rekursinya menjadi $f(n + x) = x(x+1)\cdots(x + n -
        1)\,f(x)$ dan $f(n) = (n-1)!$, turunkanlah dari pertanyaan 14
        bahwa untuk $x \in \intoc01$ berlaku
        \[
        f(x) = \lim_{n\to\infty}
        \frac{n!\,n^x}{x(x+1)\cdots(x+n)} :
        \]
        sehingga $f$ tunggal, jadi $f = \Gamma$, dan \emph{rumus
        limit Gauss} berlaku (perluaslah ke setiap $x > 0$ lewat
        rekursinya).
  \item Definisikan \emph{fungsi Beta} $B(x, y) =
        \int_0^1t^{x-1}(1-t)^{y-1}\,\dd t$ (dengan $x, y > 0$). Buktikanlah
        kekonvergenannya, rekursi $B(x+1, y) =
        \frac{x}{x+y}\,B(x, y)$ (integralkan secara parsial), dan
        $B(1, y) = \frac1y$.
  \item Tunjukkan bahwa $x \mapsto B(x, y)$ bersifat log-cembung (lewat Hölder
        lagi), lalu terapkan Bohr--Mollerup pada
        \[
        f(x) = \frac{B(x, y)\,\Gamma(x + y)}{\Gamma(y)}
        \]
        untuk menyimpulkan \emph{rumus Euler}: $B(x, y) =
        \dfrac{\Gamma(x)\Gamma(y)}{\Gamma(x+y)}$ --- tanpa integral
        rangkap di mana pun.
  \item Hitunglah $B(\frac12, \frac12)$ secara langsung (substitusikan $t
        = \sin^2\theta$) lalu turunkan $\Gamma(\frac12) =
        \sqrt\pi$: yakni integral Gauss pada Bagian I, yang dipulihkan
        oleh kecembungan belaka. Bandingkan kedua buktinya dalam satu
        kalimat masing-masing.
  \item (Penggandaan Legendre) Tunjukkan bahwa
        \[
        g(x) = \frac{2^{x-1}}{\sqrt\pi}\,
        \Gamma\Bigl(\frac x2\Bigr)
        \Gamma\Bigl(\frac{x+1}2\Bigr)
        \]
        memenuhi ketiga hipotesis Bohr--Mollerup, lalu
        simpulkan $g = \Gamma$, yakni $\Gamma(2z) =
        \frac{2^{2z-1}}{\sqrt\pi}\,\Gamma(z)\,\Gamma(z +
        \tfrac12)$ untuk setiap $z > 0$.
  \item Turunkan bentuk tertutup $\Gamma\bigl(n + \tfrac12\bigr)
        = \dfrac{(2n)!}{4^n\,n!}\sqrt\pi$, lalu buktikan, lewat
        lema kemiringan yang diterapkan pada $\log\Gamma$ di sekitar bilangan bulat
        yang besar, keasimtotikan
        \[
        \frac{\Gamma(n + \frac12)}{\Gamma(n)\,\sqrt n}
        \longrightarrow 1 .
        \]
  \item Gabungkan kedua pertanyaan terakhir menjadi keasimtotikan
        binomial pusat
        \[
        \binom{2n}{n} \sim \frac{4^n}{\sqrt{\pi n}},
        \]
        lalu periksalah secara numerik untuk $n = 10$ (dengan $\binom{20}{10} =
        184756$, terhadap $4^{10}/\sqrt{10\pi} \approx
        187079$: yakni nisbah $\approx 0.988$).
  \item (Sintesis) Konstanta $\sqrt\pi$ kini telah muncul sebagai
        integral Gauss (Bagian I), sebagai $B(\frac12,
        \frac12)$ (pertanyaan 18), dan di dalam penggandaannya
        (pertanyaan 19); sedangkan taksiran binomial pusatnya
        mendahului baik Stirling (yakni soal akhir pekan
        \cref{ch:b3:product}) maupun de Moivre--Laplace. Petakanlah
        hubungannya: pernyataan mana setara dengan yang mana,
        dan apa yang disumbangkan masing-masing teknik --- pendiferensialan di bawah
        integral lawan kecembungan --- yang tak dapat disumbangkan yang lain?
\end{enumerate}

\medskip
\noindent\textbf{Bagian V --- Tiga buah hasil lagi.}
\begin{enumerate}[resume]
  \item (Wallis, lewat Beta) Untuk $p > -1$, substitusikan $t =
        \sin^2\theta$ untuk menunjukkan bahwa
        \[
        W_p = \int_0^{\pi/2}\sin^p\theta\,\dd\theta
        = \frac12\,B\Bigl(\frac{p+1}2, \frac12\Bigr),
        \]
        lalu turunkan dari rekursi Betanya (pertanyaan 16) bahwa
        $W_{n+2} = \frac{n+1}{n+2}\,W_n$ untuk bilangan bulat $n \geq
        0$. Hitunglah $W_{2n}$ dan $W_{2n+1}$ dalam bentuk tertutup,
        tunjukkan $W_{2n+1}/W_{2n} \to 1$ lewat penjepitan, lalu simpulkan
        dengan \emph{hasil kali Wallis}
        \[
        \frac\pi2 = \lim_{n\to\infty}
        \prod_{k=1}^{n}\frac{4k^2}{4k^2 - 1} .
        \]
  \item (Fungsi Gauss berjumpa sebuah frekuensi) Untuk $b \in \R$ tetapkan
        \[
        \Phi(b) = \int_{-\infty}^{+\infty}
        \eu^{-x^2}\cos(2bx)\,\dd x .
        \]
        Tunjukkan bahwa $\Phi$ bersifat $\mathcal C^1$ pada $\R$, bahwa sebuah
        pengintegralan parsial melahirkan persamaan diferensial
        $\Phi'(b) = -2b\,\Phi(b)$, lalu simpulkan
        \[
        \Phi(b) = \sqrt\pi\,\eu^{-b^2} :
        \]
        sehingga fungsi Gauss mereproduksi dirinya di bawah transformasi ini ---
        yakni satu-satunya identitas yang akan menjalankan
        \cref{ch:b3:fouriertransform}.
  \item (Integral Frullani) Untuk $0 < a < b$, tunjukkan bahwa
        \[
        \int_0^{+\infty}
        \frac{\eu^{-ax} - \eu^{-bx}}{x}\,\dd x
        = \log\frac ba ,
        \]
        dengan menurunkannya terhadap parameter $a$ (benarkanlah
        dominasinya pada setiap $\intco{a_0}{+\infty}$, dengan $a_0 > 0$,
        lalu kenali konstantanya dengan membiarkan $a \to b$).
        Di manakah tepatnya integrannya memerlukan
        kesingularan terhapuskannya di $x = 0$?

\end{enumerate}
\end{problem}
